\makeatletter
\def\rap@fontpath{../../../Common/}
\makeatother
\documentclass{../../../Common/Latex/Templates/rapport}
% ---------------------------------------------------------------------------
% Shared settings for the H2026 General Assembly document set.
%
% Usage, directly after \documentclass{../../../Common/Latex/Templates/rapport}
% in a document one folder below this file:
%   % ---------------------------------------------------------------------------
% Shared settings for the H2026 General Assembly document set.
%
% Usage, directly after \documentclass{../../../Common/Latex/Templates/rapport}
% in a document one folder below this file:
%   % ---------------------------------------------------------------------------
% Shared settings for the H2026 General Assembly document set.
%
% Usage, directly after \documentclass{../../../Common/Latex/Templates/rapport}
% in a document one folder below this file:
%   \input{../ga-common}
%
% Image paths are relative to the compiled document's folder, so every
% document using this file must sit one folder below it.
% ---------------------------------------------------------------------------

% Meeting details, used on covers and in each document's ID block.
\newcommand{\meetingdate}{30 September 2026}
\newcommand{\meetingtime}{12:00}
\newcommand{\meetinglocation}{Festsalen, NMBU}
\newcommand{\meetingline}{\textbf{Meeting:} \meetingdate, \meetingtime, \meetinglocation}

% Header: moon logo on the left, full-text logo on the right. \projectname is
% still set because the cover title uses it.
\projectname{Eclipse NMBU}
\projectlogo{../../../Common/Eclipse Images/EclipseMoon}
\orglogo{../../../Common/Eclipse Images/EclipseFullText}

%
% Image paths are relative to the compiled document's folder, so every
% document using this file must sit one folder below it.
% ---------------------------------------------------------------------------

% Meeting details, used on covers and in each document's ID block.
\newcommand{\meetingdate}{30 September 2026}
\newcommand{\meetingtime}{12:00}
\newcommand{\meetinglocation}{Festsalen, NMBU}
\newcommand{\meetingline}{\textbf{Meeting:} \meetingdate, \meetingtime, \meetinglocation}

% Header: moon logo on the left, full-text logo on the right. \projectname is
% still set because the cover title uses it.
\projectname{Eclipse NMBU}
\projectlogo{../../../Common/Eclipse Images/EclipseMoon}
\orglogo{../../../Common/Eclipse Images/EclipseFullText}

%
% Image paths are relative to the compiled document's folder, so every
% document using this file must sit one folder below it.
% ---------------------------------------------------------------------------

% Meeting details, used on covers and in each document's ID block.
\newcommand{\meetingdate}{30 September 2026}
\newcommand{\meetingtime}{12:00}
\newcommand{\meetinglocation}{Festsalen, NMBU}
\newcommand{\meetingline}{\textbf{Meeting:} \meetingdate, \meetingtime, \meetinglocation}

% Header: moon logo on the left, full-text logo on the right. \projectname is
% still set because the cover title uses it.
\projectname{Eclipse NMBU}
\projectlogo{../../../Common/Eclipse Images/EclipseMoon}
\orglogo{../../../Common/Eclipse Images/EclipseFullText}

\usepackage{enumitem}
\usepackage{array}
\usepackage{tabularx}
\usepackage{lineno}
\usepackage{hyperref}
\hypersetup{colorlinks=true, linkcolor=blue}

\reporttitle{Case 5: Bylaw Amendment}

\begin{document}

\reportmainmatter
\begingroup
\small
\setlength{\parindent}{0pt}
\setlength{\tabcolsep}{6pt}
\renewcommand{\arraystretch}{1.2}

\begin{center}
  {\Large\bfseries Bylaw Amendment (vedtektsendring): Proposed New Bylaws}\par
  \vspace{0.15cm}
  {\normalsize 26H-ECL-GA-CASE4A \textperiodcentered{} Case 5 \textperiodcentered{} Proposed by: the Board}\par
  \vspace{0.1cm}
  {\small \meetingline}\par
\end{center}

\vspace{0.3cm}
\begin{tabularx}{\textwidth}{|>{\bfseries\raggedright\arraybackslash}p{4.2cm}|X|}
  \hline
  C-suite deliberation & The C-suite has reviewed the proposed bylaw amendment and recommends it to the General Assembly for approval. \\
  \hline
  Background           & The Old Bylaws (13 sections) no longer describe how Eclipse NMBU operates. The Board proposes a full restructuring of the Bylaws (Old Bylaws \S\S1--13) into the attached text (18 sections, five parts). The main changes are summarised in the plain-language explanation, 26H-ECL-GA-CASE4B. \\
  \hline
  Proposal             & The General Assembly resolves to replace the Old Bylaws with the attached Proposed New Bylaws, with any Amendments adopted under the voting procedure in 26H-ECL-GA-CASE4C. Adoption requires 2/3 of all members of Eclipse NMBU to vote in favour (Old Bylaws \S4.6). Individual Amendments to the text pass by simple majority. The New Bylaws take effect when the Meeting Protocol (26H-ECL-GA-PROT1) is signed. \\
  \hline
  External information & The attached text is an English translation of the Norwegian draft (``Vedtekter for Eclipse NMBU: Omstrukturert utkast, 21 Sept 2026''). The Norwegian text is the controlling version unless the General Assembly adopts the English version as controlling (proposed \S17.3). Section numbering and cross-references (jf. \S\ldots) follow the Norwegian draft, so the two can be checked against each other clause by clause. Attachments: Proposed New Bylaws (below), plain-language explanation (26H-ECL-GA-CASE4B), voting procedure (26H-ECL-GA-CASE4C), Norwegian original of the Old Bylaws (26H-ECL-GA-CASE4D) and of this proposed text (26H-ECL-GA-CASE4E). \\
  \hline
\end{tabularx}
\endgroup

\newpage
\begin{center}
  {\Large\bfseries Attachment: Proposed New Bylaws (nye vedtekter), full text}\par
\end{center}
\tableofcontents

\vsection{How to read this document}
Terminology follows the Multilingual Organisational Glossary, 26H-ECL-GA-GLOSS1: established terms are written \textbf{English (Norwegian)} on first use per section. Cross-references of the form ``jf. \S X'' point to another section of \emph{these proposed New Bylaws}, not to the Old Bylaws, unless stated otherwise; every such cross-reference below is a clickable link to the paragraph it cites. The number in the margin next to each line is that line's number in this document, for pinpointing Amendments and discussion to an exact line.

\linenumbers

\vsection{Part I: Foundations (Del I: Grunnlag)}

\vsubsection{\S 1 Name and founding (Navn og stiftelse)}\label{bylaw:1}
\textbf{\S 1.1 Name (Navn).}\label{bylaw:1.1} The Association's (foreningens) name is Eclipse NMBU.

\textbf{\S 1.2 Founding and affiliation (Stiftelse og tilknytning).}\label{bylaw:1.2} Eclipse NMBU is a voluntary student association (frivillig studentforening) at the Norwegian University of Life Sciences (NMBU), founded 9 October 2024.

\vsubsection{\S 2 Purpose and core principles (Formål og grunnprinsipper)}\label{bylaw:2}
\textbf{\S 2.1 Purpose (Formål).}\label{bylaw:2.1} Eclipse NMBU shall be a social and professional arena for students interested in space. The Association is primarily for students who are semester-registered at NMBU.

\textbf{\S 2.2 Activity (Virksomhet).}\label{bylaw:2.2} The Association's activity is voluntary and directed at research and development (R\&D). The overarching vision is to develop sustainable space technology building on the UN Sustainable Development Goals.

\textbf{\S 2.3 Core principles (Grunnprinsipper).}\label{bylaw:2.3} The Association's activity is built on the following core principles:
\begin{enumerate}
  \item \textbf{Documentation:} What is not documented in writing is not considered done.
  \item \textbf{Extreme ownership:} If something goes wrong, responsibility lies with both the leadership and the individual member.
\end{enumerate}

\textbf{\S 2.4 Ethics and environment (Etikk og miljø).}\label{bylaw:2.4} The Association shall, so far as possible, consider environmental and ethical factors when planning its activities. In dealings with other organisations, associations, events and clients, the Association shall seek out alternatives that safeguard ethical and environmental considerations.

\textbf{\S 2.5 Organisational sustainability (Organisatorisk bærekraft).}\label{bylaw:2.5} The Association's structure, build-up and organisation shall ensure that future generations of members can build on earlier activity, including through documentation, archiving and experience transfer (erfaringsoverføring), cf. \S\S10.2 and 16.5.

\vsubsection{\S 3 Organisational form (Organisasjonsform)}\label{bylaw:3}
\textbf{\S 3.1 Independence (Uavhengighet).}\label{bylaw:3.1} The Association is politically and religiously independent.

\textbf{\S 3.2 Legal status (Juridisk status).}\label{bylaw:3.2} The Association is a freestanding legal person with members, and is self-owning (selveiende). Neither members nor others have any claim on the Association's assets or property, nor are they liable for the Association's debts or other obligations.

\textbf{\S 3.3 Operating language (Driftsspråk).}\label{bylaw:3.3}
\begin{enumerate}
  \item The Association's operating language (driftsspråk) is English.
  \item All texts the Association produces in Norwegian shall be translated to English.
  \item Texts in English shall be translated to Norwegian if requested.
\end{enumerate}

\vsubsection{\S 4 Definitions (Definisjoner og ordforklaringer)}\label{bylaw:4}
\textbf{\S 4.1 Scope (Virkeområde).}\label{bylaw:4.1} The terms in this section apply throughout these Bylaws. The terms shall be used consistently in resolutions, protocols and other organisational documentation.

\textbf{\S 4.2 Organs (Organer).}\label{bylaw:4.2}
\begin{enumerate}
  \item General Assembly (GF) (generalforsamling): the Association's supreme organ, cf. \hyperref[bylaw:8]{\S 8}.
  \item Board (styret): the Association's supreme organ between General Assemblies, cf. \hyperref[bylaw:12]{\S 12}.
  \item Election and Control Committee (valgog kontrollkomité): a committee elected by the General Assembly that runs elections to offices and monitors compliance with the Bylaws, cf. \hyperref[bylaw:13]{\S 13}.
\end{enumerate}

\textbf{\S 4.3 Roles (Roller).}\label{bylaw:4.3}
\begin{enumerate}
  \item CEO (Chief Executive Officer): Chair of the Board (Styreleder).
  \item CTO (Chief Technical Officer): Technical Lead (Teknisk ansvarlig) and Deputy Chair (nestleder).
  \item CBO (Chief Business Officer): Business and Finance Lead (Forretningsog økonomiansvarlig).
  \item CPO (Chief People Officer): People and Membership Lead (Personalog medlemsansvarlig).
  \item Project Lead (Prosjektleder): a member with overall responsibility for carrying out a project.
  \item Team Lead (Teamleder): a member responsible for a team, i.e. a subject-matter group spanning multiple projects, cf. \hyperref[bylaw:12.9]{\S 12.9}.
  \item Nearest Leader (Nærmeste leder): the role a member reports to in the Association's governance structure. For an ordinary member, the nearest leader is the Team Lead. The Team Lead reports to the Project Lead of the relevant project, and the Project Lead to the responsible Board member, cf. \hyperref[bylaw:16.2]{\S 16.2}.
  \item Responsibility Role (Ansvarsverv): an assigned and delimited leadership responsibility beyond ordinary membership, including Project Lead and Team Lead.
\end{enumerate}

\textbf{\S 4.4 Membership (Medlemskap).}\label{bylaw:4.4}
\begin{enumerate}
  \item Active Member (Aktivt medlem): synonymous with ``full member''. The term Active Member is used in these Bylaws, cf. \hyperref[bylaw:5]{\S 5}.
  \item Mentor (Mentor): membership tier for people on the Association's mentor list, cf. \hyperref[bylaw:5.4]{\S 5.4}.
  \item Alumni (Alumni): membership tier for people whose Membership Contract ended without renewal, cf. \hyperref[bylaw:5.4]{\S 5.4}.
  \item Active Period (Aktiv periode): NMBU's spring and autumn parallel, from the start of term to the beginning of the examination period.
\end{enumerate}

\textbf{\S 4.5 General Assembly and elections (Generalforsamling og valg).}\label{bylaw:4.5}
\begin{enumerate}
  \item Elective General Assembly (Valggeneralforsamling): the General Assembly in the spring semester, at which offices (tillitsverv) are elected for the coming term, cf. \hyperref[bylaw:8.2]{\S 8.2}.
  \item Supplementary General Assembly (Suppleringsgeneralforsamling): the General Assembly in the autumn semester, at which supplementary elections are held for vacant offices and for the Election and Control Committee, cf. \hyperref[bylaw:8.2]{\S 8.2}.
  \item Office (Tillitsverv): an office elected by the General Assembly for a term, including Board offices and offices on the Election and Control Committee, cf. \hyperref[bylaw:10]{\S 10}.
  \item Right to Propose (Forslagsrett): the right to submit Cases and Amendments at the General Assembly, cf. \hyperref[bylaw:5.5]{\S 5.5}.
  \item Amendment (Endringsforslag): a proposal to amend a Case on the Agenda, cf. \hyperref[bylaw:9.5]{\S 9.5}.
  \item Advance Nomination (Forhåndsbenking): nomination of a candidate before the General Assembly, cf. \hyperref[bylaw:11.3]{\S 11.3}.
  \item Floor Nomination (Benkeforslag): nomination of a candidate from the floor during the General Assembly, cf. \hyperref[bylaw:11.3]{\S 11.3}.
  \item Floor Speech (Benketale): a speech in which the person who made a Floor Nomination explains the nomination.
  \item Advance Vote (Forhåndsstemme): a vote in an election, cast before voting in that election opens at the General Assembly, cf. \hyperref[bylaw:11.4]{\S 11.4}.
\end{enumerate}

\textbf{\S 4.6 Meeting officers and meeting rules (Møtefunksjonærer og møteregler).}\label{bylaw:4.6}
\begin{enumerate}
  \item Meeting Chair (Ordstyrer): chairs the General Assembly, cf. \hyperref[bylaw:9.2]{\S 9.2}.
  \item Secretary (Referent): keeps the protocol of the General Assembly, cf. \hyperref[bylaw:9.8]{\S 9.8}.
  \item Vote Counters (Tellekorps): the members who count votes at the General Assembly, cf. \hyperref[bylaw:9.3]{\S 9.3}.
  \item Protocol Signer (Protokollunderskriver): a member who checks and signs the protocol, cf. \hyperref[bylaw:9.9]{\S 9.9}.
  \item Speakers' List (Taleliste): the record of who has asked for the floor, and in what order.
  \item Speech (Innlegg): a speech in discussion of a Case, given in the order of the Speakers' List.
  \item Reply (Replikk): a short response to the immediately preceding Speech, cf. \hyperref[bylaw:9.4]{\S 9.4}.
  \item Point of Order (Innlegg til forretningsorden): a speech about the conduct of the meeting, which takes priority over the Speakers' List.
  \item Closing the Speakers' List (Strek på talelisten): a resolution that no new speakers may add themselves to the Case.
  \item Voting Order (Voteringsorden): the order in which the proposals are put to a vote, cf. \hyperref[bylaw:9.6]{\S 9.6}.
  \item Protocol Note (Protokolltilførsel): a short remark a member requires entered into the protocol, cf. \hyperref[bylaw:9.8]{\S 9.8}.
\end{enumerate}

\textbf{\S 4.7 Authority and mandate (Myndighet og fullmakt).}\label{bylaw:4.7}
\begin{enumerate}
  \item Authority (Myndighet): the decision-making power an organ or role holds directly under these Bylaws.
  \item Mandate (Fullmakt): the right to make decisions and bind the Association on the CEO's behalf, within an assigned area of responsibility, cf. \hyperref[bylaw:14]{\S 14}.
  \item Delegation (Delegasjon): transfer of Mandate from the CEO to another Board member, cf. \hyperref[bylaw:14.2]{\S 14.2}.
  \item Signature Right (Signaturrett): the right to sign binding agreements on the Association's behalf, cf. \hyperref[bylaw:14.7]{\S 14.7}.
  \item Standing Order (Føringsorden): rules adopted by the Board as an addendum to the Bylaws, cf. \hyperref[bylaw:17.4]{\S 17.4}.
\end{enumerate}

\textbf{\S 4.8 Finance and projects (Økonomi og prosjekter).}\label{bylaw:4.8}
\begin{enumerate}
  \item Free Funds (Frie midler): funds the Association disposes of without earmarking, cf. \hyperref[bylaw:15.6]{\S 15.6}.
  \item NASA Project Life Cycle: the framework for phase division and gate reviews in projects, in a modified version set out in the Association's governance documents, cf. \hyperref[bylaw:16.3]{\S 16.3}.
  \item MCR (Mission Concept Review): the concept review required to start a project, cf. \hyperref[bylaw:16.1]{\S 16.1}.
  \item SRR (System Requirements Review): the requirements review required to start a project, cf. \hyperref[bylaw:16.1]{\S 16.1}.
\end{enumerate}

\vsection{Part II: Membership (Del II: Medlemskap)}

\vsubsection{\S 5 Membership (Medlemskap)}\label{bylaw:5}
\textbf{\S 5.1 Membership tiers (Medlemsnivåer).}\label{bylaw:5.1} Eclipse NMBU operates with six membership tiers. The tiers are hierarchical, such that higher tiers hold all rights and duties belonging to the tiers below:
\begin{enumerate}
  \item \textbf{Board Member (Styremedlem):} holds the legal and administrative responsibility for the Association.
  \item \textbf{Core Member (Kjernemedlem):} a member assigned a Responsibility Role by the Board or a Board member, cf. \hyperref[bylaw:4.3]{\S 4.3} no.\ 8, including Project Lead, Team Lead or another delimited leadership responsibility.
  \item \textbf{Active Member (Aktivt medlem):} a member admitted under \S5.2.
  \item \textbf{Mentor (Mentor):} a person on the Association's mentor list, cf. \hyperref[bylaw:5.4]{\S 5.4}.
  \item \textbf{Inactive Member (Inaktivt medlem):} a member registered in the Association who is not an Active Member.
  \item \textbf{Alumni (Alumni):} a member with Alumni status, cf. \hyperref[bylaw:5.4]{\S 5.4}.
\end{enumerate}
The CPO decides whether a member is active or inactive.

\textbf{\S 5.2 Admission and Membership Contract (Opptak og medlemskontrakt).}\label{bylaw:5.2}
\begin{enumerate}
  \item All members shall sign a Membership Contract confirming that they have read the Bylaws, and that they understand and accept the risk associated with spaceflight activity.
  \item To become an Active Member, a member must have completed an interview and signed the Membership Contract.
  \item The CPO is responsible for the admission process. The interview is conducted by two Board members, who decide on admission.
\end{enumerate}

\textbf{\S 5.3 Membership period (Medlemsperiode).}\label{bylaw:5.3} The membership period follows the term of office in \S10.1 and runs from 1 August to 31 July. A person may join at any time during the period, but membership ends 31 July. Mentors and Alumni are exempt from this; their membership lasts until the person concerned withdraws, or the CPO removes them, cf. \hyperref[bylaw:5.4]{\S 5.4}.

\textbf{\S 5.4 Mentors and Alumni (Mentorer og alumni).}\label{bylaw:5.4}
\begin{enumerate}
  \item Mentor is its own membership tier, which may also be held by people outside the Association. People may be asked to become a mentor, and those who agree are entered on the Association's mentor list. The CPO is responsible for the mentor list.
  \item A member whose Membership Contract ends without being renewed automatically receives Alumni status.
\end{enumerate}

\textbf{\S 5.5 Rights at the General Assembly (Rettigheter på generalforsamlingen).}\label{bylaw:5.5}
\begin{enumerate}
  \item Board Members, Core Members and Active Members have voting rights and are eligible for office.
  \item Voting rights presuppose that the member has been admitted as an Active Member, cf. \S5.2, at least one month before the General Assembly.
  \item Members with voting rights, and Mentors, have Right to Propose.
  \item All members, including Inactive Members, Mentors and Alumni, have meeting and speaking rights.
  \item Inactive Members, Mentors and Alumni have neither voting rights nor eligibility for office.
  \item The General Assembly may, by simple majority, grant meeting and speaking rights to people who are not members.
\end{enumerate}

\textbf{\S 5.6 Duties (Plikter).}\label{bylaw:5.6}
\begin{enumerate}
  \item Members are bound by resolutions made by the General Assembly and the Board.
  \item \textbf{Duty to deliver:} a member who has taken on a task is obliged to deliver it within the set time frame.
  \item \textbf{Duty to notify:} if a member cannot carry out a task taken on, the member shall, without undue delay, notify their Nearest Leader, cf. \hyperref[bylaw:4.3]{\S 4.3} no.\ 7. The Nearest Leader is responsible for the task being reassigned.
  \item \textbf{Absence:} if a member is inactive without valid reason for a continuous period of four weeks in an Active Period, the Board may consider ending the membership. The reason for absence shall be clarified with the member's Nearest Leader.
\end{enumerate}

\textbf{\S 5.7 Board authority in membership matters (Styrets myndighet i medlemssaker).}\label{bylaw:5.7} Outside the General Assembly, the Board may, by majority resolution:
\begin{enumerate}
  \item Change a member's tier on the grounds of a failure to meet duties, except that assessment of whether a member is active or inactive is excepted, cf. \hyperref[bylaw:5.1]{\S 5.1}.
  \item Grant or refuse applications for offices and membership, and change who holds an office.
  \item Remove members from the Association's membership list if \S5.6 no.\ 4 is breached.
\end{enumerate}

\vsubsection{\S 6 Whistleblowing and reactions (Varsling og reaksjoner)}\label{bylaw:6}
\textbf{\S 6.1 Definition (Definisjon).}\label{bylaw:6.1} Whistleblowing (varsling) is when a member, a Board member, or a person in contact with the Association, reports a reprehensible matter (kritikkverdig forhold) connected with the Association's activity.

\textbf{\S 6.2 Reprehensible matters (Kritikkverdige forhold).}\label{bylaw:6.2} Reprehensible matters include, among others:
\begin{enumerate}
  \item Matters relating to health, environment and safety.
  \item Personal conflicts.
  \item Disruptive conduct.
  \item Legal violations, including theft, financial breach of trust and embezzlement.
  \item Sexual harassment.
  \item Assault.
\end{enumerate}

\textbf{\S 6.3 Case handling (Saksbehandling).}\label{bylaw:6.3}
\begin{enumerate}
  \item A report is directed to the CPO, or to the CEO if the report concerns the CPO.
  \item In the event of serious allegations of a legal violation, the matter must be reported to the police before the Board may process the report.
  \item In cases directed at one or more persons, all involved shall be given the opportunity to account for their version of the situation.
  \item Legal violations and sexual harassment are not accepted, and the Board may report matters to the police based on its own observations.
\end{enumerate}

\textbf{\S 6.4 Impartiality (Habilitet).}\label{bylaw:6.4} When handling reports, the Board's impartiality shall be assessed. A Board member who is not impartial does not take part in the handling. If the Board as a whole is not impartial, an independent body shall be appointed to handle the matter.

\textbf{\S 6.5 Reactions (Reaksjoner).}\label{bylaw:6.5} A resolution on a reaction requires a 2/3 majority on the Board, cf. \hyperref[bylaw:12.7]{\S 12.7} no.\ 5. The following reactions may be decided:
\begin{enumerate}
  \item Changes to the Association's routines.
  \item A warning.
  \item Suspension or exclusion from office and the Association's activities for a period.
  \item Permanent exclusion from office and the Association's activities.
\end{enumerate}

\vsubsection{\S 7 Privacy (Personvern)}\label{bylaw:7}
\textbf{\S 7.1 Processing of personal data (Behandling av personopplysninger).}\label{bylaw:7.1} The Association processes only personal data that is necessary and relevant to administering membership. Processing shall follow the Personal Data Act (personopplysningsloven) and the guidelines of the Norwegian Data Protection Authority (Datatilsynet).

\textbf{\S 7.2 Retention on end of membership (Oppbevaring ved opphør av medlemskap).}\label{bylaw:7.2} On end of membership, the Association retains the member's personal data until a request for deletion is made under \S7.3.

\textbf{\S 7.3 Request for deletion (Anmodning om sletting).}\label{bylaw:7.3}
\begin{enumerate}
  \item A member or former member may at any time request that their own personal data be deleted.
  \item The request shall be made in writing to the CPO.
  \item Deletion shall be carried out without undue delay, and the member shall receive confirmation that deletion has been carried out.
\end{enumerate}

\textbf{\S 7.4 Duty to inform (Informasjonsplikt).}\label{bylaw:7.4} The Association shall, on admission and on end of membership, inform the member of the right to request deletion under \S7.3.


\vsection{Part III: Organs (Del III: Organer)}

\vsubsection{\S 8 General Assembly (Generalforsamling)}\label{bylaw:8}
\textbf{\S 8.1 Supreme organ (Øverste organ).}\label{bylaw:8.1} The General Assembly is the Association's supreme organ.

\textbf{\S 8.2 Frequency and types (Hyppighet og typer).}\label{bylaw:8.2}
\begin{enumerate}
  \item The General Assembly shall be held at least once per semester, within the end of November in the autumn semester and within the end of April in the spring semester.
  \item The General Assembly in the spring semester is the Elective General Assembly; the General Assembly in the autumn semester is the Supplementary General Assembly.
\end{enumerate}

\textbf{\S 8.3 Notice and deadlines (Innkalling og frister).}\label{bylaw:8.3}
\begin{enumerate}
  \item Notice shall reach members at least one month before the General Assembly.
  \item Cases to be handled by the General Assembly, including proposed Bylaw Amendments, shall reach the Board at least two weeks before the General Assembly.
  \item Case Papers shall be available to members at least one week before the General Assembly.
\end{enumerate}

\textbf{\S 8.4 Quorum and majority requirements (Vedtaksførhet og flertallskrav).}\label{bylaw:8.4}
\begin{enumerate}
  \item The General Assembly is quorate when lawfully convened, irrespective of the number attending.
  \item The General Assembly passes resolutions by simple majority unless these Bylaws provide otherwise. Majority requirements for Bylaw Amendments and dissolution follow \S\S17.1 and 18.1.
  \item No member has more than one vote, and voting may not be by proxy. Advance Voting in elections is governed by \S11.4.
\end{enumerate}

\textbf{\S 8.5 Extraordinary General Assembly (Ekstraordinær generalforsamling).}\label{bylaw:8.5} An extraordinary General Assembly may be called if at least 1/3 of members require it, or if the Board finds it necessary. \S\S8.3 and 8.4 apply correspondingly. Only the Cases stated in the notice are handled at an extraordinary General Assembly.

\textbf{\S 8.6 Agenda (Saksliste).}\label{bylaw:8.6} The Agendas of the Elective and Supplementary General Assemblies follow the order below.

\emph{Elective General Assembly (spring semester):}
\begin{enumerate}
  \item Constitution: (1) election of Meeting Chair; (2) election of Secretary; (3) election of Vote Counters; (4) election of Protocol Signers; (5) approval of Notice; (6) approval of Agenda.
  \item Presentation and discussion of Cases: (1) Semester Report; (2) previous semester's financial statement; (3) coming semester's budget; (4) submitted Cases.
  \item Voting on Cases.
  \item Election of CEO, CTO, CBO and CPO for the coming term, cf. \hyperref[bylaw:11]{\S 11}.
  \item Other business.
\end{enumerate}

\emph{Supplementary General Assembly (autumn semester):}
\begin{enumerate}
  \item Constitution: (1) election of Meeting Chair; (2) election of Secretary; (3) election of Vote Counters; (4) election of Protocol Signers; (5) approval of Notice; (6) approval of Agenda.
  \item Presentation and discussion of Cases: (1) Semester Report; (2) previous semester's financial statement; (3) coming semester's budget; (4) submitted Cases.
  \item Voting on Cases.
  \item Supplementary elections, cf. \hyperref[bylaw:11]{\S 11}: (1) supplementary elections to vacant offices for the remainder of the term; (2) election of a chair and two members to the Election and Control Committee.
  \item Other business.
\end{enumerate}

Each Case under point 2 is presented and discussed together, and the Cases are voted on under point 3. Case handling follows \S9.

\vsubsection{\S 9 Meeting rules for the General Assembly (Møteregler for generalforsamlingen)}\label{bylaw:9}
\textbf{\S 9.1 Opening and constitution (Åpning og konstituering).}\label{bylaw:9.1}
\begin{enumerate}
  \item The General Assembly is opened by the CEO, or by the CTO under \S12.3. Whoever opens the meeting leads the election of the Meeting Chair.
  \item Until the Vote Counters are elected, votes are counted by whoever leads the meeting.
  \item Once the Meeting Chair is elected, the Meeting Chair takes over the conduct of the meeting and carries out the remaining points of the constitution.
  \item When approving the Notice, the General Assembly considers whether the Notice and Case Papers comply with \S8.3. If the Notice is not approved, the General Assembly cannot pass resolutions.
\end{enumerate}

\textbf{\S 9.2 Meeting Chair (Ordstyrer).}\label{bylaw:9.2}
\begin{enumerate}
  \item The Meeting Chair leads the General Assembly in accordance with the Bylaws and the approved Agenda, and shall ensure orderly and neutral conduct of the meeting.
  \item The Meeting Chair keeps the Speakers' List, gives the floor, and may take the floor from a speaker if the speaking time is exceeded or the speech does not concern the Case.
  \item The Meeting Chair may set speaking time for Speeches and Replies.
  \item The Meeting Chair decides questions on the conduct of the meeting. The decision may be referred to the General Assembly, which decides the question by simple majority.
  \item The Meeting Chair does not take part in debate. If the Meeting Chair wishes to take part in debate on a Case, conduct of the meeting for that Case is handed to a person elected by the General Assembly.
  \item The Meeting Chair should not be a Board member.
\end{enumerate}

\textbf{\S 9.3 Vote Counters (Tellekorps).}\label{bylaw:9.3}
\begin{enumerate}
  \item The Vote Counters consist of two members, who shall be neutral in the Cases they count.
  \item The Vote Counters register the number of eligible voters present at the constitution, and keep the number updated through the meeting.
  \item The Vote Counters distribute and collect ballots, verify that only Eligible Voters cast a vote, count the votes, and report the result to the Meeting Chair.
  \item The Vote Counters count votes on Cases, and in elections not conducted by the Election and Control Committee, cf. \hyperref[bylaw:11.5]{\S 11.5}.
  \item Ballots are kept by the Vote Counters until the protocol is signed, and are shredded thereafter.
\end{enumerate}

\textbf{\S 9.4 Presentation and discussion of Cases (Presentasjon og diskusjon av saker).}\label{bylaw:9.4}
\begin{enumerate}
  \item Each Case is presented by the proposer or the Board, and is discussed before the next Case is presented.
  \item During discussion, the floor is given in the order requested.
  \item After a Speech, a Reply may be given. Whoever gave the Speech has the right to a Reply.
  \item A Point of Order takes priority over the Speakers' List.
  \item The Meeting Chair, or any Eligible Voter, may propose Closing the Speakers' List. The General Assembly decides the proposal by simple majority.
\end{enumerate}

\textbf{\S 9.5 Amendments (Endringsforslag).}\label{bylaw:9.5} An Amendment to a Case on the Agenda must be submitted in writing before voting on Cases begins.

\textbf{\S 9.6 Voting on Cases (Stemming over saker).}\label{bylaw:9.6}
\begin{enumerate}
  \item Before voting, the Meeting Chair puts forward a Voting Order, which is approved by the General Assembly.
  \item Amendments are voted on before the original proposal. If there are several Amendments to the same point, the most far-reaching proposal is voted on first.
  \item Voting is by show of hands, unless at least one Eligible Voter requires written and secret voting.
  \item Once voting is underway, the floor is given only for questions about the vote itself.
  \item Blank votes are counted as not cast.
  \item In the event of a tied vote, the proposal is considered not adopted.
\end{enumerate}

\textbf{\S 9.7 Other business (Eventuelt).}\label{bylaw:9.7}
\begin{enumerate}
  \item Matters for other business must be submitted in writing before voting on Cases begins.
  \item Under other business, only matters submitted under no.\ 1 are handled.
  \item No resolution on Bylaw Amendments, elections, financial dispositions or dissolution may be made under other business.
\end{enumerate}

\textbf{\S 9.8 Secretary and protocol (Referent og protokoll).}\label{bylaw:9.8}
\begin{enumerate}
  \item The Secretary keeps the protocol of the General Assembly.
  \item The protocol shall at least contain: (1) time and place of the meeting; (2) number of eligible voters present; (3) who was elected during the constitution; (4) all proposals made, with the proposer; (5) resolutions with vote counts; (6) the result of all elections; (7) Protocol Notes.
  \item Any Eligible Voter may require a short Protocol Note entered into the protocol.
  \item The protocol shall be signed by the Protocol Signers and made available to members no later than two weeks after the General Assembly.
\end{enumerate}

\textbf{\S 9.9 Protocol Signers (Protokollunderskrivere).}\label{bylaw:9.9}
\begin{enumerate}
  \item The General Assembly elects two Protocol Signers from among the members present.
  \item The Protocol Signers check that the protocol gives a correct account of the General Assembly, and sign it. They may require corrections before signing.
\end{enumerate}

\vsubsection{\S 10 Offices (Tillitsverv)}\label{bylaw:10}
\textbf{\S 10.1 Term of office (Vervperiode).}\label{bylaw:10.1} All offices in the Association, including all Board offices, run from 1 August to 31 July of the following year. The Election and Control Committee is exempt from this, cf. \hyperref[bylaw:13.1]{\S 13.1}.

\textbf{\S 10.2 Overlap and experience transfer (Overlapp og erfaringsoverføring).}\label{bylaw:10.2} The period from the Elective General Assembly until 1 August is an overlap period in which the incoming and outgoing office holders act simultaneously. In this period, the incoming office holder has decision-making authority and voting rights, while the outgoing office holder has an advisory function.

The outgoing office holder shall, in the overlap period, ensure full experience transfer, including handover of documentation, access, ongoing matters and points of contact, and shall be available to the Board and their successor.

\textbf{\S 10.3 Vacant offices (Ledige verv).}\label{bylaw:10.3}
\begin{enumerate}
  \item If an office holder steps down before the term of office has ended, the Board may, by simple majority, appoint a temporary member to the office until it is filled at an ordinary General Assembly.
  \item If an office is not filled at the General Assembly, the Board may appoint a member to the office until the next General Assembly.
  \item Nos.\ 1 and 2 apply correspondingly to the Election and Control Committee. If nobody is elected to the Committee, the Board may appoint a temporary committee until the next General Assembly.
\end{enumerate}

\vsubsection{\S 11 Elections (Valg)}\label{bylaw:11}
\textbf{\S 11.1 Scope (Virkeområde).}\label{bylaw:11.1} This section applies to all elections held at the General Assembly, including election of Meeting Chair, Secretary, Vote Counters, Protocol Signers, Board members and members of the Election and Control Committee.

\textbf{\S 11.2 Candidates (Kandidater).}\label{bylaw:11.2}
\begin{enumerate}
  \item Voting rights and eligibility for office follow \S5.5.
  \item A candidate must have consented to the candidature. Candidates who cannot attend may stand with a written candidature.
\end{enumerate}

\textbf{\S 11.3 Advance Nomination and Floor Nomination (Forhåndsbenking og benkeforslag).}\label{bylaw:11.3}
\begin{enumerate}
  \item Before the General Assembly, a call for Advance Nominations is sent out. The Election and Control Committee asks those proposed whether they wish to stand. Those who wish to stand are listed as candidates on the Agenda. The deadline for Advance Nomination is the same as for submitting Cases, cf. \hyperref[bylaw:8.3]{\S 8.3}.
  \item During the General Assembly, Floor Nomination may only be made during presentation and discussion of Cases, cf. \hyperref[bylaw:8.6]{\S 8.6}.
  \item Whoever makes a Floor Nomination is entitled to give a Floor Speech.
\end{enumerate}

\textbf{\S 11.4 Advance Voting (Forhåndsstemming).}\label{bylaw:11.4}
\begin{enumerate}
  \item Active Members with voting rights may cast an Advance Vote in elections to offices. Advance Voting does not apply to other Cases.
  \item Advance Votes are cast digitally through a solution that ensures only Eligible Voters can vote, that each member casts one vote, and that the vote is secret.
  \item The deadline for Advance Voting is when voting opens in the relevant election at the General Assembly.
  \item A member who has cast an Advance Vote may not vote again in the same election.
  \item The Election and Control Committee administers Advance Votes and counts them together with the votes cast at the General Assembly.
  \item An Advance Vote applies to the candidate it was cast for, including if new candidates are added by Floor Nomination at the General Assembly.
  \item In a second round of voting, the Advance Vote is counted if the candidate it was cast for takes part in the round. Otherwise it is not counted.
\end{enumerate}

\textbf{\S 11.5 Conduct (Gjennomføring).}\label{bylaw:11.5}
\begin{enumerate}
  \item For election to offices, the Election and Control Committee leads the election and presents the office and the candidates. Other elections are led by the Meeting Chair, except election of the Meeting Chair itself, cf. \hyperref[bylaw:9.1]{\S 9.1}.
  \item Candidates are given the opportunity to present themselves, and the assembly the opportunity to ask questions.
  \item The election is conducted in writing when there is more than one candidate for an office, or when at least one Eligible Voter requires it. Otherwise the election may be conducted by show of hands.
  \item For election to offices, the Election and Control Committee counts the votes and reports the result to the Meeting Chair. For other elections, votes are counted per \S\S9.1 and 9.3. The result is entered in the protocol.
\end{enumerate}

\textbf{\S 11.6 Voting rounds (Avstemningsrunder).}\label{bylaw:11.6}
\begin{enumerate}
  \item The first voting round is conducted with a blank vote as an alternative.
  \item The candidate who receives the most votes is elected.
  \item If the blank-vote alternative receives the most votes, a second voting round is conducted between the two candidates with the most votes, without a blank-vote alternative.
  \item If two or more candidates are tied for the most votes, a second voting round is conducted between those candidates, without a blank-vote alternative.
  \item If there is only one candidate for the office, the candidate is also elected if the blank-vote alternative receives the most votes.
  \item If no candidate is elected in the second voting round, it is repeated until a candidate is elected.
\end{enumerate}

\vsubsection{\S 12 The Board (Styret)}\label{bylaw:12}
\textbf{\S 12.1 Composition (Sammensetning).}\label{bylaw:12.1} The Board consists of four members:
\begin{enumerate}
  \item CEO: Chief Executive Officer
  \item CTO: Chief Technical Officer
  \item CBO: Chief Business Officer
  \item CPO: Chief People Officer
\end{enumerate}
All Board members are elected by the General Assembly, cf. \hyperref[bylaw:1]{\S 1}1, for the term of office in \S10.1.

\textbf{\S 12.2 CEO.}\label{bylaw:12.2} The CEO has overall administrative leadership of the Association and represents the Association externally. The CEO has the ultimate responsibility for the Association operating within Norwegian law and the University's regulations, and for strategic planning, organisational culture, HSE protocols and coordination between the other Board roles. The CEO's mandate and duty to delegate follow \S14.

\textbf{\S 12.3 CTO.}\label{bylaw:12.3} The CTO has the overall responsibility for the Association's technical projects, including system integration, technical safety and technical progress. The CTO shall ensure that technical development follows set timelines and milestones, cf. \hyperref[bylaw:1]{\S 1}6, and is responsible for resource management and logistics relating to technical equipment and workshop space. The CTO calls technical status meetings.

The CTO is the Association's Deputy Chair (nestleder) and takes over the CEO's responsibility when the CEO is unavailable. Beyond this, the CTO has no greater authority than the other Board members.

\textbf{\S 12.4 CBO.}\label{bylaw:12.4} The CBO is responsible for the Association's economy, cf. \hyperref[bylaw:15]{\S 15}. This includes financial planning, bookkeeping, budget control, auditing, insurance and financial risk management. The CBO further has responsibility for sponsor relations, partnership agreements, branding and external communication, including content and maintenance of the Association's digital platforms.

\textbf{\S 12.5 CPO.}\label{bylaw:12.5} The CPO is responsible for the Association's members and internal organisation, including recruitment, the admission process, member follow-up and administration of membership tiers, cf. \hyperref[bylaw:5]{\S 5}. The CPO further has responsibility for internal communication, organisational culture and follow-up of HSE, in cooperation with the CEO, and for case handling of reports, cf. \hyperref[bylaw:6]{\S 6}.

\textbf{\S 12.6 Board meetings (Styremøter).}\label{bylaw:12.6} The Board has quorum when at least three of the Board's members, including constituted members, are present. Resolutions are made by ordinary majority of the votes cast, unless these Bylaws provide otherwise. All Board meetings shall be recorded in a protocol.

\textbf{\S 12.7 Voting on the Board (Stemmegivning i styret).}\label{bylaw:12.7}
\begin{enumerate}
  \item Each Board member has one vote. The CEO's vote counts as 1.5 votes on all Board matters.
  \item When the CTO has taken over the CEO's responsibility under \S12.3, the CTO's vote counts as 1.5 votes. The additional 0.5 weighting cannot fall to other Board members.
  \item A constituted Board member has voting rights, cf. \hyperref[bylaw:10.3]{\S 10.3}.
  \item In the event of a tied vote, the proposal is considered not adopted.
  \item Where these Bylaws require a 2/3 majority on the Board, the majority is calculated from the total voting weight of the Board members taking part in the vote.
\end{enumerate}

\textbf{\S 12.8 Duty of loyalty (Lojalitetsplikt).}\label{bylaw:12.8} The Board's members, including the CEO, shall in Board matters put the interests of the Association as a whole ahead of other considerations.

\textbf{\S 12.9 Teams (Team).}\label{bylaw:12.9}
\begin{enumerate}
  \item Teams are established by the Board member responsible for the subject area, and span multiple projects.
  \item The Board member appoints a Team Lead for the team.
\end{enumerate}

\vsubsection{\S 13 The Election and Control Committee (Valgog kontrollkomiteen)}\label{bylaw:13}
\textbf{\S 13.1 Composition and term (Sammensetning og vervperiode).}\label{bylaw:13.1}
\begin{enumerate}
  \item The Election and Control Committee consists of three members: one chair and two members.
  \item The chair and members are elected at the Supplementary General Assembly in the autumn semester.
  \item The term of office runs from 1 January following the Supplementary General Assembly to 31 December of the same year.
  \item Board members may not be members of the Committee.
\end{enumerate}

\textbf{\S 13.2 Tasks (Oppgaver).}\label{bylaw:13.2}
\begin{enumerate}
  \item The Committee shall find suitable candidates for election, administer Advance Votes, and conduct elections to offices, cf. \hyperref[bylaw:11]{\S 11}.
  \item The Committee shall monitor compliance with the Bylaws, and shall report to the General Assembly.
\end{enumerate}

\textbf{\S 13.3 Duty of confidentiality and eligibility (Taushetsplikt og valgbarhet).}\label{bylaw:13.3}
\begin{enumerate}
  \item The Committee's members are bound by a duty of confidentiality.
  \item Members of the Committee may not stand for election.
\end{enumerate}


\vsubsection{\S 14 Mandate, delegation and Signature Right (Fullmakt, delegasjon og signaturrett)}\label{bylaw:14}
\textbf{\S 14.1 Starting point (Utgangspunkt).}\label{bylaw:14.1} The CEO holds the Association's overall Mandate within the limits set by the General Assembly and the Board, and may enter binding agreements on the Association's behalf within those limits, cf. \S\S14.7 and 14.8.

\textbf{\S 14.2 Duty to delegate (Delegasjonsplikt).}\label{bylaw:14.2} The CEO shall delegate Mandate to the other Board members within their respective areas of responsibility, cf. \S\S12.3--12.5. The delegation shall be presented to, and approved by, the General Assembly. Authority over the Association's economy belongs to the CBO directly, cf. \hyperref[bylaw:1]{\S 1}5.2, and is not covered by the delegation. Delegation of Signature Right follows \S14.9.

\textbf{\S 14.3 Form and content (Form og innhold).}\label{bylaw:14.3} The delegation shall be in writing and shall state which areas of responsibility, categories of decision and any financial limits it covers. The delegation applies for the term of office, cf. \hyperref[bylaw:1]{\S 1}0.1, unless otherwise decided, and shall be recorded in the General Assembly's protocol.

\textbf{\S 14.4 Responsibility (Ansvar).}\label{bylaw:14.4} Delegation does not relieve the CEO of the overall responsibility for the Association's activity.

\textbf{\S 14.5 Change and withdrawal (Endring og tilbakekalling).}\label{bylaw:14.5} Change or withdrawal of delegated Mandate is handled by the General Assembly. Where circumstances cannot wait for an ordinary General Assembly, the Board may, by 2/3 majority, temporarily suspend a delegated Mandate until the next General Assembly.

\textbf{\S 14.6 Onward delegation (Videredelegering).}\label{bylaw:14.6} A Board member may onward-delegate delimited parts of their Mandate to a Project Lead or Team Lead for specific tasks. Onward delegation shall be in writing and notified to the Board, and does not relieve the Board member of responsibility.

\textbf{\S 14.7 Signature Right (Signaturrett).}\label{bylaw:14.7} The CEO holds the Association's Signature Right.

\textbf{\S 14.8 Conditions for signing (Vilkår for signering).}\label{bylaw:14.8} Signing a binding agreement presupposes a prior Board resolution.

\textbf{\S 14.9 Delegation of Signature Right (Delegasjon av signaturrett).}\label{bylaw:14.9} The CEO may delegate Signature Right directly to the other Board members, without General Assembly handling. The delegation shall be documented in writing and does not relieve the CEO of responsibility, cf. \hyperref[bylaw:14.4]{\S 14.4}.

\textbf{\S 14.10 Responsibility for contracts (Ansvar for kontrakter).}\label{bylaw:14.10} The sitting Board is responsible for all agreements in force.

\textbf{\S 14.11 Digital storage (Digital oppbevaring).}\label{bylaw:14.11} All documents of the Association, including agreements and contracts signed on paper, shall also be stored digitally. A document signed on paper shall be scanned or otherwise reproduced in digital form without undue delay; the physical signature does not exempt the document from this requirement.

\vsection{Part IV: Operations (Del IV: Drift)}

\vsubsection{\S 15 Economy (Økonomi)}\label{bylaw:15}
\textbf{\S 15.1 Use of funds (Bruk av midler).}\label{bylaw:15.1} The Association's funds shall be used only in accordance with the Association's purpose, cf. \hyperref[bylaw:2]{\S 2}.

\textbf{\S 15.2 Authority over the economy (Myndighet over økonomien).}\label{bylaw:15.2} The CBO has full authority over the Association's economy. This authority follows directly from these Bylaws and is not delegated from the CEO, cf. \hyperref[bylaw:14.2]{\S 14.2}.

\textbf{\S 15.3 Budget and accounts (Budsjett og regnskap).}\label{bylaw:15.3} A budget and accounts for the Association's funds, with supporting vouchers, shall always exist. The budget and accounts are reviewed and approved by the General Assembly, cf. \hyperref[bylaw:8.6]{\S 8.6}.

\textbf{\S 15.4 Approval of payments (Godkjenning av utbetalinger).}\label{bylaw:15.4} Payments are approved by the CBO.

\textbf{\S 15.5 Financial investments (Økonomiske investeringer).}\label{bylaw:15.5} Resolutions on financial investments require a 2/3 majority on the Board, cf. \hyperref[bylaw:12.7]{\S 12.7} no.\ 5.

\textbf{\S 15.6 Free Funds (Frie midler).}\label{bylaw:15.6}
\begin{enumerate}
  \item Funds the Association earns that are not earmarked on receipt shall be placed in a dedicated account for Free Funds.
  \item Members may apply for support from Free Funds for projects. The application is made to the CBO.
  \item The CBO assesses applications and decides which are forwarded to the Board. The Board resolves on allocation.
\end{enumerate}

\textbf{\S 15.7 Reporting (Rapportering).}\label{bylaw:15.7} The CBO shall follow up all financial matters and keep the Board continuously informed of the Association's financial dispositions.

\vsubsection{\S 16 Project execution (Prosjektgjennomføring)}\label{bylaw:16}
\textbf{\S 16.1 Starting a project (Oppstart av prosjekt).}\label{bylaw:16.1} All larger projects shall be approved by the Board before start-up. To start a project, an MCR and an SRR shall be delivered to the Board. The Board assesses the project against the Association's purpose, resource situation and feasibility, and approves or rejects it at the meeting where the MCR and SRR are presented.

\textbf{\S 16.2 Project organisation (Prosjektorganisering).}\label{bylaw:16.2} The Board appoints a Project Lead for each project. The Project Lead reports to the CTO in technical projects, and otherwise to the Board member responsible for the subject area.

\textbf{\S 16.3 Project model (Prosjektmodell).}\label{bylaw:16.3} Projects are carried out according to a modified version of the NASA Project Life Cycle. Phase division, gate reviews and documentation requirements are set out in the Association's governance documents.

\textbf{\S 16.4 Deviations (Avvik).}\label{bylaw:16.4} Material deviations from an approved plan, including changes to requirements, budget, timeline or safety-relevant assumptions, shall be presented to the Board.

\textbf{\S 16.5 Closure and documentation (Avslutning og dokumentasjon).}\label{bylaw:16.5} A project is closed with a final report to the Board. Project documentation shall be archived so that later members can build on the work, cf. \hyperref[bylaw:2.5]{\S 2.5}.

\vsection{Part V: Governance Documents and Dissolution (Del V: Styringsdokumenter og opphør)}

\vsubsection{\S 17 Bylaw Amendments and governance documents (Vedtektsendringer og styringsdokumenter)}\label{bylaw:17}
\textbf{\S 17.1 Bylaw Amendments (Vedtektsendringer).}\label{bylaw:17.1} Bylaw Amendments may only happen when (1) at least 2/3 of the Association's members are present at the General Assembly at the time of the vote, and (2) at least 2/3 of the members present vote in favour of the amendment. A proposed Bylaw Amendment must be submitted within the deadline in \S8.3.

\textbf{\S 17.2 Entry into force (Ikrafttredelse).}\label{bylaw:17.2} Bylaw Amendments take effect when the protocol of the General Assembly is signed, unless the General Assembly resolves otherwise.

\textbf{\S 17.3 Language versions (Språkversjoner).}\label{bylaw:17.3} In the event of conflict between language versions of the Bylaws, the version adopted by the General Assembly controls, cf. \hyperref[bylaw:3.3]{\S 3.3}.

\textbf{\S 17.4 Standing Orders (Føringsorden).}\label{bylaw:17.4} The Board may adopt a Standing Order as an addendum to the Bylaws. A Standing Order may not conflict with the Bylaws, and lapses on a change of Board.

\vsubsection{\S 18 Dissolution (Opphør)}\label{bylaw:18}
\textbf{\S 18.1 Dissolution (Oppløsning).}\label{bylaw:18.1} The Association may be dissolved if 2/3 of the eligible voting members at the General Assembly vote in favour of dissolution. A Case on dissolution shall be submitted within the deadline in \S8.3.

\textbf{\S 18.2 Disposal of assets (Disponering av eiendeler).}\label{bylaw:18.2} On dissolution, the Association's funds and property are transferred to another suitable student association with the same purpose, chosen by the General Assembly. No member has a claim on the Association's funds or any share of them, cf. \hyperref[bylaw:3.2]{\S 3.2}.

\end{document}
